% =========================================================================
\section{Presión electrostática}\label{sec:presion}
% =========================================================================
Si el conductor está cargado, su carga reside en la superficie y cada elemento de carga es repelido por todos los demás. Aparece así una
fuerza por unidad de área que tiende a expulsar la carga hacia el exterior: la \emph{presión electrostática}. Calcularla exige un poco de
cuidado, porque el campo en la superficie es discontinuo: vale $\sigma/\eo$ justo fuera y $0$ justo dentro. ¿Cuál de los dos actúa sobre la carga?

% -------------------------------------------------------------------------
\subsection{Un argumento ingenuo y su fallo}\label{sec:presion-ingenuo}
% -------------------------------------------------------------------------
Una primera idea sería multiplicar la carga de un elemento $\Delta S$ por el campo exterior:
\[
  \Pel\overset{?}{=}\frac{F}{\Delta S}=\frac{\Delta q\,E}{\Delta S}=\frac{\sigma\,\Delta S\,(\sigma/\eo)}{\Delta S}=\frac{\sigma^{2}}{\eo}.
\]
El razonamiento tiene un fallo: el campo $\sigma/\eo$ incluye el que crea la propia carga $\Delta q$, y \emph{una distribución de carga no ejerce
fuerza neta sobre sí misma} (por la tercera ley de Newton, las fuerzas internas se cancelan dos a dos), igual que nadie puede levantarse a sí
mismo tirando de sus propias manos hacia arriba. Hay que separar el campo en la contribución del elemento y la del resto de la carga.

% -------------------------------------------------------------------------
\subsection{El campo que actúa sobre la carga superficial}\label{sec:presion-campo}
% -------------------------------------------------------------------------
% =========================================================================
%  Figura: presión electrostática. (a) Elemento DeltaS de la superficie, campo
%  del disco y campo del resto en dos puntos 1 (fuera) y 2 (dentro).
%  (b) Componente normal de cada contribución frente a la coordenada xi
%  normal a la superficie: E_disco salta de -sigma/2eps0 a +sigma/2eps0,
%  E_resto es continuo y la suma pasa de 0 a sigma/eps0.
% =========================================================================
\begin{figure}[t]
\centering
\begin{tikzpicture}[font=\small]
  % ---------------------------------------------------------------- (a)
  \begin{scope}
    % superficie curvada del conductor (vista en perspectiva)
    \fill[top color=azulvelo!70!white, bottom color=metal!80!white, draw=azulprof, line width=0.6pt]
      (-3.0,-0.3) .. controls (-1.6,0.55) and (0.9,0.65) .. (3.0,0.05)
      -- (2.25,-1.05) .. controls (0.3,-0.55) and (-1.9,-0.75) .. (-3.6,-1.2) -- cycle;
    \node[font=\footnotesize, text=azulprof] at (-2.55,-0.95) {$S'$};
    % disco DeltaS
    \fill[cargapos!25] (0,-0.12) ellipse (0.62 and 0.17);
    \draw[cargapos, line width=0.6pt] (0,-0.12) ellipse (0.62 and 0.17);
    \node[font=\footnotesize] at (0.92,-0.32) {$\Delta S$};
    % campo del disco (a ambos lados) y del resto
    \draw[-{Stealth[length=4pt]}, verdeazul, line width=0.9pt] (0.12,-0.05) -- (0.12,1.15);
    \draw[-{Stealth[length=4pt]}, verdeazul, line width=0.9pt, densely dashed] (0.12,-0.2) -- (0.12,-1.4);
    \node[font=\footnotesize, text=verdeazul] at (0.72,0.95) {$\Edisc$};
    \draw[-{Stealth[length=4pt]}, cargapos, line width=0.9pt] (-0.22,-0.9) -- (-0.22,1.55);
    \node[font=\footnotesize, text=cargapos] at (-0.78,1.42) {$\Erest$};
    \draw[-{Stealth[length=3.5pt]}, azulprof, line width=0.6pt] (1.75,0.35) -- (1.75,0.95) node[right, font=\footnotesize] {$\un$};
    % puntos 1 y 2
    \begin{scope}[shift={(2.85,-1.95)}]
      \draw[azulacero, line width=1.3pt] (-0.55,0) -- (0.55,0);
      \fill[azulprof] (0,0.22) circle (1.4pt) node[right=3pt, font=\footnotesize] {1 (fuera)};
      \fill[azulprof] (0,-0.22) circle (1.4pt) node[right=3pt, font=\footnotesize] {2 (dentro)};
      \node[font=\footnotesize, anchor=east] at (-0.6,0) {$\Delta S$};
    \end{scope}
    \draw[->, grismedio, line width=0.4pt] (0.65,-0.18) .. controls (1.6,-0.7) .. (2.25,-1.85);
    \node at (0,-2.95) {(a)};
  \end{scope}
  % ---------------------------------------------------------------- (b)
  \begin{scope}[xshift=6.3cm, yshift=-2.1cm]
    \begin{axis}[tema, width=7.0cm, height=5.4cm, axis lines=left,
        xlabel={$\xi$ (normal a la superficie)}, ylabel={$E_n$}, xmin=-2.1, xmax=2.1, ymin=-0.75, ymax=1.25,
        xtick={0}, xticklabels={$0$}, ytick={-0.5,0,0.5,1},
        yticklabels={$-\frac{\sigma}{2\eo}$,$0$,$\frac{\sigma}{2\eo}$,$\frac{\sigma}{\eo}$},
        grid=none, legend columns=1,
        legend style={at={(0.5,1.04)}, anchor=south, font=\scriptsize}]
      \addplot[verdeazul, line width=1.1pt, densely dashed, const plot] coordinates {(-2,-0.5) (0,0.5) (2,0.5)};
      \addlegendentry{campo del disco $E_{\mathrm{disco}}$}
      \addplot[cargapos, line width=1.1pt] coordinates {(-2,0.52) (2,0.52)};
      \addlegendentry{campo del resto $E_{\mathrm{resto}}$}
      \addplot[azulprof, line width=1.4pt, const plot] coordinates {(-2,0.0) (0,1) (2,1)};
      \addlegendentry{campo total $E=E_{\mathrm{disco}}+E_{\mathrm{resto}}$}
      \draw[punteado] (axis cs:0,-0.75) -- (axis cs:0,1.25);
      \node[font=\scriptsize, text=grispizarra] at (axis cs:-1.05,-0.63) {conductor};
      \node[font=\scriptsize, text=grispizarra] at (axis cs:1.05,-0.63) {vacío};
    \end{axis}
    \node at (2.6,-1.25) {(b)};
  \end{scope}
\end{tikzpicture}
\caption{Presión electrostática. (a) El campo en un elemento $\Delta S$ de la superficie es la suma del campo del propio disco, discontinuo
($\pm\sigma/2\eo$ a cada lado), y del campo $\Erest$ del resto de la carga, continuo. (b) Componente normal de cada contribución frente a la
coordenada $\xi$ perpendicular a la superficie: $\Erest$ duplica el campo fuera y lo anula dentro, y es el único que actúa sobre la carga de
$\Delta S$.}
\label{fig:presion}
\end{figure}


\begin{lema}[Campo del resto de la carga]\label{lem:resto}
Sea $\Delta S$ un elemento de la superficie de un conductor en equilibrio, tan pequeño que $\sigma$ es constante en él, y sea $\Erest$ el campo
creado en $\Delta S$ por toda la carga del sistema \emph{excepto} la de $\Delta S$. Entonces
\begin{equation}\label{eq:campo-resto}
  \resalta{\Erest=\frac12\bigl(\vE_1+\vE_2\bigr)=\frac{\sigma}{2\eo}\,\un ,}
\end{equation}
donde $\vE_1$ y $\vE_2$ son los campos totales justo fuera y justo dentro del conductor, respectivamente.
\end{lema}
\begin{proof}
Descomponemos el campo total en dos puntos muy próximos a $\Delta S$, uno fuera (1) y otro dentro (2), como suma del campo del disco $\Delta S$ y
del campo del resto (\figpanel{fig:presion}{a}): $\vE=\Edisc+\Erest$.

\emph{El campo del resto es continuo} en $\Delta S$: es el campo que crea en un punto del vacío una distribución de carga (la de $S'-\Delta S$) que
no tiene carga en ese punto, de modo que no presenta ninguna discontinuidad allí; toma el mismo valor $\Erest$ en 1 y en 2.

\emph{El campo del disco es discontinuo.} En puntos suficientemente próximos a $\Delta S$, el disco se ve como un plano indefinido con densidad
$\sigma$, cuyo campo es $\sigma/(2\eo)$, perpendicular al plano y dirigido hacia fuera de él a ambos lados~\cite{TEMA1}:
\[
  \vE_{\mathrm{disco},1}=+\frac{\sigma}{2\eo}\,\un,\qquad\vE_{\mathrm{disco},2}=-\frac{\sigma}{2\eo}\,\un .
\]
Por tanto, usando \cref{eq:campo-superficie} y \cref{eq:campo-nulo},
\begin{alignat*}{2}
  \vE_1&=\Erest+\frac{\sigma}{2\eo}\,\un=\frac{\sigma}{\eo}\,\un &\qquad&\text{(punto 1, fuera)},\\
  \vE_2&=\Erest-\frac{\sigma}{2\eo}\,\un=\vc{0} &&\text{(punto 2, dentro)}.
\end{alignat*}
Sumando miembro a miembro ambas igualdades, $\vE_1+\vE_2=2\,\Erest$, de donde
\[
  \Erest=\frac12\bigl(\vE_1+\vE_2\bigr)=\frac12\Bigl(\frac{\sigma}{\eo}\,\un+\vc{0}\Bigr)=\frac{\sigma}{2\eo}\,\un .
\]
\end{proof}

\begin{observacion}[Lectura del resultado]\label{obs:resto}
El campo del resto de la carga vale exactamente la mitad del campo exterior, y es el \emph{promedio} de los campos a ambos lados de la superficie.
Su papel es doble (\figpanel{fig:presion}{b}): fuera, se suma al del disco y \emph{duplica} la intensidad hasta $\sigma/\eo$; dentro, se opone a él y lo
\emph{cancela} por completo. El equilibrio del conductor es, desde este punto de vista, un ajuste fino entre la carga de cada elemento y la de todos
los demás.
\end{observacion}

% -------------------------------------------------------------------------
\subsection{La presión electrostática}\label{sec:presion-valor}
% -------------------------------------------------------------------------
\begin{teorema}[Presión electrostática]\label{teo:presion}
Sobre la superficie de un conductor en equilibrio actúa una fuerza por unidad de área
\begin{equation}\label{eq:presion}
  \resalta{\Pel=\frac{\dif F}{\dif S}=\frac{\sigma^{2}}{2\eo}=\frac{\eo}{2}\,E^{2},}\qquad
  \frac{\dif\vF}{\dif S}=\frac{\sigma^{2}}{2\eo}\,\un ,
\end{equation}
dirigida siempre hacia el exterior del conductor, cualquiera que sea el signo de $\sigma$; aquí $E=\abs{\sigma}/\eo$ es el campo justo fuera de la superficie.
\end{teorema}
\begin{proof}
La fuerza sobre la carga $\Delta q=\sigma\,\Delta S$ del elemento la ejerce solo el resto de la carga, de modo que, por el \cref{lem:resto},
\[
  \Delta\vF=\Delta q\,\Erest=\sigma\,\Delta S\,\frac{\sigma}{2\eo}\,\un
  \quad\Longrightarrow\quad\Pel=\frac{\Delta F}{\Delta S}=\frac{\sigma^{2}}{2\eo} .
\]
Como $\sigma^{2}\ge0$, la fuerza tiene siempre el sentido de $\un$, hacia fuera. Por último, $\sigma^{2}/(2\eo)=\tfrac{\eo}{2}(\sigma/\eo)^{2}=\tfrac{\eo}{2}E^{2}$.
\end{proof}

\begin{observacion}[Presión y densidad de energía]\label{obs:presion-energia}
La expresión $\Pel=\tfrac{\eo}{2}E^{2}$ coincide con la densidad de energía del campo electrostático, $u=\tfrac{\eo}{2}E^{2}$, que estudiaremos más
adelante. No es casual: si un elemento de superficie se desplaza hacia fuera una distancia $\dif\xi$, el volumen $\dif S\,\dif\xi$, en el que antes había
un campo $E$, pasa a formar parte del metal y queda sin campo; el trabajo de la fuerza electrostática, $\Pel\,\dif S\,\dif\xi$, es igual a la energía
$u\,\dif S\,\dif\xi$ que desaparece.
\end{observacion}

\begin{ejemplo}[Fuerza entre los dos hemisferios de una esfera cargada]\label{ej:hemisferios}
Una esfera conductora de radio $R$ tiene carga $Q$. ¿Con qué fuerza se repelen sus dos hemisferios?

\emph{Resolución.} La carga se reparte uniformemente, $\sigma=Q/(4\pi R^{2})$, y la presión $\Pel=\sigma^{2}/(2\eo)$ es la misma en toda la superficie. Por
simetría, la fuerza sobre el hemisferio norte apunta según $\uz$, y solo contribuye la componente $\un\cdot\uz=\cos\theta$:
\begin{align*}
  F_z&=\iint_{\text{hemisferio}}\Pel\cos\theta\,\dsc
     =\frac{\sigma^{2}}{2\eo}\int_{0}^{2\pi}\!\!\int_{0}^{\pi/2}\cos\theta\,R^{2}\sen\theta\,\dif\theta\,\dif\varphi\\
     &=\frac{\sigma^{2}}{2\eo}\,2\pi R^{2}\,\frac12=\frac{\sigma^{2}\pi R^{2}}{2\eo}=\frac{Q^{2}}{32\pi\eo R^{2}} .
\end{align*}
El resultado tiene una lectura geométrica: es la presión por el área de la \emph{proyección} del hemisferio sobre el plano ecuatorial, $\pi R^{2}$.
\end{ejemplo}

\begin{observacion}[Órdenes de magnitud]\label{obs:presion-orden}
El campo máximo que soporta el aire seco antes de ionizarse es $\Erup\simeq3\cdot10^{6}\,\mathrm{V/m}$ (\cref{sem:globo}); a ese campo,
$\Pel=\tfrac{\eo}{2}\Erup^{2}\simeq40\,\mathrm{Pa}$, unas $2500$ veces menor que la presión atmosférica. La presión electrostática es pequeña,
pero perceptible en membranas ligeras: una pompa de jabón cargada se hincha ligeramente.
\end{observacion}
